% secV5_change.tex -- Section 7. CONTAINER (27 Oct 2026): text to be written by the author.
% Word budget: 450 words (JPART count: running text, footnotes and citations as printed).
% Outline, numbers and sources: referee_EE_2026-10/ESQUELETO.md, section 7. The archived text
% (_arquivo_texto_IA_2026-10-27/) is not a source.

\section{The timing of requests and referrals}\label{sec:service}

An unchanged definition of the waiting list does not imply an unchanged composition. Within the window of complete coverage established in Section~\ref{sec:coverage}, children began entering the request record at younger ages, while referrals within a year of the first request became more frequent. These are descriptive changes in the recorded flows, not estimates of policy effects.

\subsection{Earlier first requests}\label{sec:timing}

Among the 2012--2016 birth cohorts, the share with a first request before age 12 months increased successively from 38.3\% to 41.5\%, 44.1\%, 48.6\%, and 50.1\% of live births (Table~\ref{tab:timing}). The corresponding shares before six and 24 months rose from 23.5\% to 33.4\% and from 59.4\% to 70.9\%, respectively. The change is not solely an increase in the fraction eventually requesting childcare: among children requesting before 24 months, the share doing so before 12 months increased from 64.4\% to 70.7\%. With fixed district weights, it increased from 64.3\% to 70.8\%, and the within-district share rose in 93 of 96 districts. The median age at first request among requests filed in a given month also declined, from 43 months in January 2013 to 13 months in December 2018 (Figure~\ref{fig:F8}). Earlier evidence documents a decline in age at first registration \citep{soares2021ecd}; the implication here concerns the composition of a dated administrative count.

An early request does not necessarily represent a reservation for a later start. Among children requesting before six months of age who received a referral within 30 days, 94.3\% subsequently enrolled, at a median age of four months (Table~\ref{tab:C4}). The desired start date is not recorded, so these data do not identify families' preferred timing of care.

\subsection{Referrals and enrolment}

The cumulative incidence of referral within 12 months of a first request, accounting for ageing out as a competing risk, rose from 63.1\% for the 2012 cohort to 87.6\% for the 2016 cohort (Table~\ref{tab:service}). This increase is positive across alternative specifications, although its magnitude ranges from 12 to 31 percentage points (Table~\ref{tab:specs}). Referral is distinct from enrolment: 9.1\% of children making requests in 2013--2017 have no recorded enrolment, while enrolment within 12 months rose from 52.3\% to 74.7\% between the 2012 and 2016 cohorts (Tables~\ref{tab:T4d_never} and~\ref{tab:service}). These comparisons do not identify the cause of faster referrals.

Earlier requests and more frequent referrals within a year jointly change which requests remain unresolved at a given date. A snapshot of the waiting list therefore combines changes in entry timing and subsequent trajectories. This motivates the request-date flow indicators developed in Section~\ref{sec:indicators}.
